\documentclass[french]{yLectureNote}

\title{Fiche de révision}
\subtitle{Synthèse des notions essentielles}
\author{Paulhenry Saux}
\date{\today}
\yLanguage{Français}

\professor{}

\usepackage{graphicx}
\usepackage[utf8]{inputenc}
\usepackage{geometry}
\usepackage{tikz}
\usepackage{tkz-tab}
\usepackage{tabularx}
\usepackage{awesomebox}
\usepackage{menukeys}
\usepackage{fancyhdr}
\usepackage{blindtext}
\usepackage{hyperref}
\usepackage{caption}
\usepackage{pifont}
\usepackage{array}
\usepackage{lipsum}
\usepackage{yFlatTable}
\usepackage{multicol}

\newcommand{\Lim}[1]{\lim\limits_{\substack{#1}}\:}
\renewcommand{\vec}{\overrightarrow}
\newcommand{\bz}{\overline{z}}
\DeclareMathOperator{\card}{card}
\renewcommand{\vec}{\overrightarrow}
\newcommand{\N}[0]{\mathbb{N}}
\newcommand{\R}[0]{\mathbb{R}}
\newcommand{\C}[0]{\mathbb{C}}
\newcommand{\dd}[0]{\mathrm{d}}
\newcommand{\er}{\vec{e_{\rho}}}
\newcommand{\ep}{\vec{e_{\varphi}}}
\newcommand{\pip}{\pi^+}
\newcommand{\pim}{\pi^-}
\newcommand{\mv}{\mathcal{V}}
\DeclareMathOperator{\Div}{Div}
\newcommand{\rot}{\vec{\mathrm{rot}}}
\newcommand{\grad}{\vec{\mathrm{grad}}}

\begin{document}
%Foche méthode : résoudre un probleme de thermo simple + pièges à éviter
\section{Système thermodynamique}

\begin{definition}
[Système thermodynamique]
    Ensemble matériel impliqué dans une processus de conversion d'énergie. 
    Il peut \^etre délimité par des frontière fixes ou mobiles réelles ou imaginaires

\end{definition}

\begin{definition}
[Système isolé]
    N'échange ni matière ni énergie avec l'extérieur du système. Contraire : système ouvert
\end{definition}

\begin{definition}
[Système fermé]
    Il n'y a pas de transfert de matière mais échange d'énergie. La masse est constante

\end{definition}

% \begin{definition}
% [Système ouvert]
%     Échange matière et énergie avec le reste de l'univers

% \end{definition}

\section{Variables}

\begin{definition}
[Variable thermodynamique]
    Toute grandeur variant, mesurable, directement ou non, appartenant au système

\end{definition}

\begin{definition}
[Densité massique d'une grandeur extensive]
    Soit Y une grandeur extensive. On définit $$y = \Lim{m\to 0}\frac{Y}{m}$$ qui est une grandeur intensive.

\end{definition}

\begin{theorem}
Pour un système fermé contenant un corps pur, l'état est connu si 2 variables thermodynamiques, souvent (P,T), (P,V), (T,V).

\end{theorem}

\section{Notion de température}

\begin{definition}
[Température]
    La grandeur que partagent 2 systèmes en équilibre thermique

\end{definition}

\begin{theorem}
[Passage d'une unité à l'autre]
    \begin{itemize}
        \item $T[F] = \frac95 T[C]+32$
        \item $T[K] = T[C]+273.15$
    \end{itemize}

\end{theorem}
\section{Pression}

\begin{definition}
[Pression]
    Notée $p$, elle s'exprime en Pascal et correspond à la densité de force qu'exerce le gaz/liquide sur son contenant :$$\dd \vec{F} = p\vec{\dd S}$$

\end{definition}

\begin{theorem}
[Loi de l'hydrostatique]
    $$\vec{\text{grad}}(P) = \rho \vec{g}$$

\end{theorem}

% \begin{theorem}
% [Loi du nivellement barométrique]
%     $$\ln\left(\frac{P(z_B)}{P(z_A)}\right) = -\frac{Mg}{RT}(z_B - z_A)$$

% \end{theorem}

\section{Équilibre Thermodynamique}

\begin{definition}
[Système à l'équilibre]
    Un système est à l'équilibre s'il n'y a pas de flux de masse, pas de flux d'énergie et que l'état est stationnaire.

\end{definition}

\section{Équation d'état}

\begin{definition}
[Gaz parfait]
    Un gaz est dit parfait s'il respecte $PV = nRT$.

\end{definition}

\begin{definition}
[Modèle de Van der Waals]
    Pour prendre en compte le volume propre des molécules ($b$) et leurs interactions ($a$), on utilise :
    $$\left( P + \frac{n^2 a}{V^2} \right) (V - nb) = nRT$$

\end{definition}

\section{Réservoir}

\begin{definition}
[Thermostat/Réservoir/Source]
    C'est un lieu à qui on peut prendre ou donner de l'énergie autant que l'on veut sans que sa température ne change.

\end{definition}

\section{Transformations}

\begin{center}
\begin{tabularx}{\textwidth}{|>{\raggedright\arraybackslash}X|>{\raggedright\arraybackslash}X|}
\hline
\textbf{Type de transformation} & \textbf{Définition / Condition} \\
\hline
Transformation & Passage d'un système d'un état initial (d'équilibre) à un autre d'équilibre aussi. \\
\hline
Transformation quasi-statique & Succession continue d'états d'équilibre, quelle que soit l'amplitude des variations. \\
\hline
Transformation non quasi-statique & Au moins un état n'est pas un état d'équilibre. \\
\hline
Transformation réversible & Peut être effectuée dans le sens inverse, retrouvant l'état initial. \\
\hline
Transformation monotherme & $T_{ext} = T_i = T_f$  \\
\hline
Transformation isobare & $P = \text{cte}$ (pression constante) \\
\hline
Transformation isochore & $V = \text{cte}$ (volume constant) \\
\hline
\end{tabularx}
\end{center}

\end{document}
